% Options for packages loaded elsewhere
\PassOptionsToPackage{unicode}{hyperref}
\PassOptionsToPackage{hyphens}{url}
\documentclass[
  polish,
  11pt,
  a4paper,
]{article}
\usepackage{xcolor}
\usepackage[margin=2cm]{geometry}
\usepackage{amsmath,amssymb}
\setcounter{secnumdepth}{-\maxdimen} % remove section numbering
\usepackage{iftex}
\ifPDFTeX
  \usepackage[T1]{fontenc}
  \usepackage[utf8]{inputenc}
  \usepackage{textcomp} % provide euro and other symbols
\else % if luatex or xetex
  \usepackage{unicode-math} % this also loads fontspec
  \defaultfontfeatures{Scale=MatchLowercase}
  \defaultfontfeatures[\rmfamily]{Ligatures=TeX,Scale=1}
\fi
\usepackage{lmodern}
\ifPDFTeX\else
  % xetex/luatex font selection
\fi
% Use upquote if available, for straight quotes in verbatim environments
\IfFileExists{upquote.sty}{\usepackage{upquote}}{}
\IfFileExists{microtype.sty}{% use microtype if available
  \usepackage[]{microtype}
  \UseMicrotypeSet[protrusion]{basicmath} % disable protrusion for tt fonts
}{}
\makeatletter
\@ifundefined{KOMAClassName}{% if non-KOMA class
  \IfFileExists{parskip.sty}{%
    \usepackage{parskip}
  }{% else
    \setlength{\parindent}{0pt}
    \setlength{\parskip}{6pt plus 2pt minus 1pt}}
}{% if KOMA class
  \KOMAoptions{parskip=half}}
\makeatother
\usepackage{color}
\usepackage{fancyvrb}
\newcommand{\VerbBar}{|}
\newcommand{\VERB}{\Verb[commandchars=\\\{\}]}
\DefineVerbatimEnvironment{Highlighting}{Verbatim}{commandchars=\\\{\}}
% Add ',fontsize=\small' for more characters per line
\usepackage{framed}
\definecolor{shadecolor}{RGB}{248,248,248}
\newenvironment{Shaded}{\begin{snugshade}}{\end{snugshade}}
\newcommand{\AlertTok}[1]{\textcolor[rgb]{0.94,0.16,0.16}{#1}}
\newcommand{\AnnotationTok}[1]{\textcolor[rgb]{0.56,0.35,0.01}{\textbf{\textit{#1}}}}
\newcommand{\AttributeTok}[1]{\textcolor[rgb]{0.13,0.29,0.53}{#1}}
\newcommand{\BaseNTok}[1]{\textcolor[rgb]{0.00,0.00,0.81}{#1}}
\newcommand{\BuiltInTok}[1]{#1}
\newcommand{\CharTok}[1]{\textcolor[rgb]{0.31,0.60,0.02}{#1}}
\newcommand{\CommentTok}[1]{\textcolor[rgb]{0.56,0.35,0.01}{\textit{#1}}}
\newcommand{\CommentVarTok}[1]{\textcolor[rgb]{0.56,0.35,0.01}{\textbf{\textit{#1}}}}
\newcommand{\ConstantTok}[1]{\textcolor[rgb]{0.56,0.35,0.01}{#1}}
\newcommand{\ControlFlowTok}[1]{\textcolor[rgb]{0.13,0.29,0.53}{\textbf{#1}}}
\newcommand{\DataTypeTok}[1]{\textcolor[rgb]{0.13,0.29,0.53}{#1}}
\newcommand{\DecValTok}[1]{\textcolor[rgb]{0.00,0.00,0.81}{#1}}
\newcommand{\DocumentationTok}[1]{\textcolor[rgb]{0.56,0.35,0.01}{\textbf{\textit{#1}}}}
\newcommand{\ErrorTok}[1]{\textcolor[rgb]{0.64,0.00,0.00}{\textbf{#1}}}
\newcommand{\ExtensionTok}[1]{#1}
\newcommand{\FloatTok}[1]{\textcolor[rgb]{0.00,0.00,0.81}{#1}}
\newcommand{\FunctionTok}[1]{\textcolor[rgb]{0.13,0.29,0.53}{\textbf{#1}}}
\newcommand{\ImportTok}[1]{#1}
\newcommand{\InformationTok}[1]{\textcolor[rgb]{0.56,0.35,0.01}{\textbf{\textit{#1}}}}
\newcommand{\KeywordTok}[1]{\textcolor[rgb]{0.13,0.29,0.53}{\textbf{#1}}}
\newcommand{\NormalTok}[1]{#1}
\newcommand{\OperatorTok}[1]{\textcolor[rgb]{0.81,0.36,0.00}{\textbf{#1}}}
\newcommand{\OtherTok}[1]{\textcolor[rgb]{0.56,0.35,0.01}{#1}}
\newcommand{\PreprocessorTok}[1]{\textcolor[rgb]{0.56,0.35,0.01}{\textit{#1}}}
\newcommand{\RegionMarkerTok}[1]{#1}
\newcommand{\SpecialCharTok}[1]{\textcolor[rgb]{0.81,0.36,0.00}{\textbf{#1}}}
\newcommand{\SpecialStringTok}[1]{\textcolor[rgb]{0.31,0.60,0.02}{#1}}
\newcommand{\StringTok}[1]{\textcolor[rgb]{0.31,0.60,0.02}{#1}}
\newcommand{\VariableTok}[1]{\textcolor[rgb]{0.00,0.00,0.00}{#1}}
\newcommand{\VerbatimStringTok}[1]{\textcolor[rgb]{0.31,0.60,0.02}{#1}}
\newcommand{\WarningTok}[1]{\textcolor[rgb]{0.56,0.35,0.01}{\textbf{\textit{#1}}}}
\usepackage{longtable,booktabs,array}
\usepackage{caption}
\captionsetup[table]{skip=6pt}
\newcounter{none} % for unnumbered tables
\usepackage{calc} % for calculating minipage widths
% Correct order of tables after \paragraph or \subparagraph
\usepackage{etoolbox}
\makeatletter
\patchcmd\longtable{\par}{\if@noskipsec\mbox{}\fi\par}{}{}
\makeatother
% Allow footnotes in longtable head/foot
\IfFileExists{footnotehyper.sty}{\usepackage{footnotehyper}}{\usepackage{footnote}}
\makesavenoteenv{longtable}
\usepackage{graphicx}
\makeatletter
\newsavebox\pandoc@box
\newcommand*\pandocbounded[1]{% scales image to fit in text height/width
  \sbox\pandoc@box{#1}%
  \Gscale@div\@tempa{\textheight}{\dimexpr\ht\pandoc@box+\dp\pandoc@box\relax}%
  \Gscale@div\@tempb{\linewidth}{\wd\pandoc@box}%
  \ifdim\@tempb\p@<\@tempa\p@\let\@tempa\@tempb\fi% select the smaller of both
  \ifdim\@tempa\p@<\p@\scalebox{\@tempa}{\usebox\pandoc@box}%
  \else\usebox{\pandoc@box}%
  \fi%
}
% Set default figure placement to htbp
\def\fps@figure{htbp}
\makeatother
\ifLuaTeX
\usepackage[bidi=basic,shorthands=off]{babel}
\else
\usepackage[bidi=default,shorthands=off]{babel}
\fi
\ifLuaTeX
  \usepackage{selnolig} % disable illegal ligatures
\fi
\setlength{\emergencystretch}{3em} % prevent overfull lines
\providecommand{\tightlist}{%
  \setlength{\itemsep}{0pt}\setlength{\parskip}{0pt}}
\usepackage{fontspec}
\defaultfontfeatures{Ligatures=TeX}
\usepackage{amsmath,amssymb}
\usepackage{booktabs}
\usepackage{array}
\usepackage{geometry}
\usepackage{enumitem}
\usepackage{xcolor}
\usepackage{tcolorbox}
\usepackage{fancyhdr}
\usepackage{titlesec}
\usepackage{multicol}

\geometry{margin=1.8cm, top=2.5cm, bottom=2cm}
\pagestyle{fancy}
\fancyhf{}
\rhead{\small Projektowanie badań — 2026/27}
\lhead{\small Zarządzanie informacją | ISI UJ}
\cfoot{\thepage}

\titleformat{\section}{\large\bfseries\color{blue!70!black}}{\thesection.}{0.5em}{}
\titleformat{\subsection}{\normalsize\bfseries}{\thesubsection.}{0.5em}{}
\titleformat{\subsubsection}{\small\bfseries\color{gray!70!black}}{\thesubsubsection.}{0.5em}{}
\titlespacing*{\section}{0pt}{1.5ex}{1ex}
\titlespacing*{\subsection}{0pt}{1ex}{0.5ex}

\setlist[itemize]{nosep, leftmargin=1.5em}
\setlist[enumerate]{nosep, leftmargin=1.5em}

\tcbset{boxrule=0.5pt, left=4pt, right=4pt, top=4pt, bottom=4pt, fonttitle=\bfseries\small}

\newtcolorbox{definicja}[1][]{colback=blue!5, colframe=blue!50!black, title={#1}}
\newtcolorbox{uwagabox}[1][]{colback=orange!10, colframe=orange!70!black, title={#1}}
\newtcolorbox{wzorbox}[1][]{colback=gray!5, colframe=gray!50!black, title={#1}}
\newtcolorbox{przykladbox}[1][]{colback=purple!5, colframe=purple!50!black, title={#1}}

\usepackage{bookmark}
\IfFileExists{xurl.sty}{\usepackage{xurl}}{} % add URL line breaks if available
\urlstyle{same}
% fallback for those not using the hyperref driver hyperxmp:
\makeatletter
\@ifundefined{xmpquote}{\newcommand{\xmpquote}[1]{#1}}{}
\makeatother
\hypersetup{
  pdftitle={Handout C03 --- Opis danych i dwa wykresy},
  pdfauthor={Marek Deja},
  pdflang={pl-PL},
  hidelinks,
  pdfcreator={LaTeX via pandoc}}

\title{Handout C03 --- Opis danych i dwa wykresy}
\usepackage{etoolbox}
\makeatletter
\providecommand{\subtitle}[1]{% add subtitle to \maketitle
  \apptocmd{\@title}{\par {\large #1 \par}}{}{}
}
\makeatother
\subtitle{Zarządzanie informacją --- część ilościowa --- 2026/27}
\author{Marek Deja}
\date{2026/27}

\begin{document}
\maketitle

\section{Plan i dane}\label{plan-i-dane}

10 min start → 20 min miary → 25 min wspólne wykresy → 25 min Z03 → 10
min oddanie. Odtwórz własny projekt przez
\texttt{rozpocznij\_zajecia("C03",\ ...)}. Dane są syntetyczne; surowego
CSV nie zmieniaj.

\section{Tabela przed wykresem}\label{tabela-przed-wykresem}

Liczba osób i liczba ważnych czasów mogą być różne. Obliczamy czas w
minutach po regułach C02. Odczytaj medianę i IQR razem ze średnią i SD.

\begin{Shaded}
\begin{Highlighting}[]
\NormalTok{dane }\OtherTok{\textless{}{-}}\NormalTok{ badaniaZI}\SpecialCharTok{::}\FunctionTok{przygotuj\_ankiete}\NormalTok{(badaniaZI}\SpecialCharTok{::}\FunctionTok{dane\_przykladowe}\NormalTok{())}\SpecialCharTok{$}\NormalTok{dane}
\NormalTok{x }\OtherTok{\textless{}{-}}\NormalTok{ dane}\SpecialCharTok{$}\NormalTok{czas\_wyszukiwania}
\FunctionTok{c}\NormalTok{(}\AttributeTok{N =} \FunctionTok{nrow}\NormalTok{(dane), }\AttributeTok{wazne =} \FunctionTok{sum}\NormalTok{(}\SpecialCharTok{!}\FunctionTok{is.na}\NormalTok{(x)), }\AttributeTok{braki =} \FunctionTok{sum}\NormalTok{(}\FunctionTok{is.na}\NormalTok{(x)),}
  \AttributeTok{srednia =} \FunctionTok{mean}\NormalTok{(x, }\AttributeTok{na.rm =} \ConstantTok{TRUE}\NormalTok{), }\AttributeTok{mediana =} \FunctionTok{median}\NormalTok{(x, }\AttributeTok{na.rm =} \ConstantTok{TRUE}\NormalTok{),}
  \AttributeTok{SD =}\NormalTok{ stats}\SpecialCharTok{::}\FunctionTok{sd}\NormalTok{(x, }\AttributeTok{na.rm =} \ConstantTok{TRUE}\NormalTok{), }\AttributeTok{IQR =}\NormalTok{ stats}\SpecialCharTok{::}\FunctionTok{IQR}\NormalTok{(x, }\AttributeTok{na.rm =} \ConstantTok{TRUE}\NormalTok{))}
\end{Highlighting}
\end{Shaded}

\begin{verbatim}
##          N      wazne      braki    srednia    mediana         SD        IQR
## 148.000000 142.000000   6.000000  10.128169   8.250000   5.854444   7.175000
\end{verbatim}

{\def\LTcaptype{none} % do not increment counter
\begin{longtable}[]{@{}ll@{}}
\toprule\noalign{}
Element & Co czytam \\
\midrule\noalign{}
\endhead
\bottomrule\noalign{}
\endlastfoot
Średnia & wrażliwa na skrajne wartości \\
Mediana & środek uporządkowanych wartości \\
SD & rozproszenie w jednostce zmiennej \\
IQR & Q3 − Q1, środkowe 50\% wartości \\
Odsetek grupy & liczebność / N ważnych grup \\
\end{longtable}
}

\section{Dwa wykresy}\label{dwa-wykresy}

Histogram: przedziały czasu, kształt/ogon; gęstość nie jest procentem w
punkcie. Pudełko: mediana, kwartyle, wąsy 1,5 IQR, punkty do
sprawdzenia. Każdy wykres ma nazwy osi, minuty, źródło syntetyczne i N
ważnych obserwacji. Punkt poza wąsami nie jest automatycznie błędem.

\section{Z03 --- wymagane wyniki}\label{z03-wymagane-wyniki}

Własny czas: N/osoby/braki, średnia, mediana, SD, IQR, minimum i
maksimum. Grupy: N i procenty z mianownikiem. Zapisz
\texttt{wyniki/opis.csv}, \texttt{histogram.png}, \texttt{boxplot.png};
dodaj linki w \texttt{odpowiedzi.md}. Własny opis z liczbami, argument
za średnią/medianą i ograniczenie do próby.

Ocena: obliczenia/mianowniki 3 pkt; dwa właściwe i czytelne wykresy 3
pkt; uzasadnienie miary 2 pkt; opis i granice wniosku 2 pkt. Pełne
poziomy: \texttt{rubryka("Z03")}.

Zapis → wykonanie skryptu od początku → \texttt{sprawdz\_zadanie("Z03")}
→ \texttt{oddaj\_zadanie("Z03")} → \texttt{status\_oddania("Z03")} →
wylogowanie R i przeglądarki.

\section{Notatki}\label{notatki}

N ważne: \ldots\ldots\ldots. Średnia: \ldots\ldots\ldots. Mediana:
\ldots\ldots\ldots. SD: \ldots\ldots\ldots. IQR: \ldots\ldots\ldots.

Co na histogramie uzasadnia wybór miary?
\ldots\ldots\ldots\ldots\ldots\ldots\ldots\ldots\ldots\ldots\ldots\ldots\ldots\ldots\ldots\ldots\ldots\ldots..

Jakie zdanie wykraczałoby poza dane?
\ldots\ldots\ldots\ldots\ldots\ldots\ldots\ldots\ldots\ldots\ldots\ldots\ldots\ldots\ldots\ldots\ldots\ldots\ldots..

\end{document}
